% Six distinctions between model levels
\begin{tikzpicture}[x=1cm,y=1cm,>=Latex]
  \tikzset{
    dist/.style={draw=blue!45!black,fill=blue!4,rounded corners,minimum width=5.0cm,minimum height=1.45cm,align=left,font=\scriptsize}
  }
  \node[dist] (r1) at (0,0) {\textbf{\#1}\quad L1 $\to$ 器件层\\无限周期布洛赫模 $\neq$ 有限器件工作模};
  \node[dist,right=0.5cm of r1] (r2) {\textbf{\#2}\quad L1 $\to$ L2\\被动冷腔共振 $\neq$ 有源阈值};
  \node[dist,below=0.45cm of r1] (r3) {\textbf{\#3}\quad L2 $\to$ L3/L4\\光学阈值 $g_{\mathrm{mod,th}},g_{\mathrm{mat,th}}$ $\neq$ 阈值电流 $I_{\mathrm{th}}$};
  \node[dist,right=0.5cm of r3] (r4) {\textbf{\#4}\quad L2 $\to$ L3/L4\\均匀增益、有效折射率代理 $\neq$ 真实注入分布};
  \node[dist,below=0.45cm of r3] (r5) {\textbf{\#5}\quad 2D $\to$ 3D\\二维教学模型 $\neq$ 三维平板器件};
  \node[dist,right=0.5cm of r5] (r6) {\textbf{\#6}\quad L2 $\to$ 阈上\\线性阈值代理 $\neq$ 阈上稳态激光理论};
\end{tikzpicture}
